\section{Caso de estudio: Modelado del proceso de secado de café}

\subsection{Planteamiento del problema}

El secado del café es una etapa crítica en la agroindustria, ya que influye directamente en la calidad final del grano, el tiempo de procesamiento y el consumo energético.

Variables como la temperatura del aire y la humedad relativa afectan la cinética de secado, el contenido final de humedad y el perfil sensorial del café. Sin embargo, estas relaciones no siempre son evidentes, lo que hace necesario el uso de herramientas estadísticas y modelos predictivos.

El objetivo de este estudio es analizar la influencia de variables de proceso sobre la cinética de secado y la calidad del café, mediante el uso de análisis estadístico y modelos de regresión.

\subsection{Variables del estudio}

Se definieron las siguientes variables:

\begin{itemize}
    \item Variables independientes:
    \begin{itemize}
        \item Temperatura del aire de secado ($^\circ$C)
        \item Humedad relativa (\%)
    \end{itemize}
    
    \item Variables de respuesta:
    \begin{itemize}
        \item Contenido de humedad del grano (\%)
        \item Tiempo de secado (h)
        \item Calidad sensorial (puntaje)
    \end{itemize}
\end{itemize}

\subsection{Diseño experimental}

Se propone un diseño factorial donde se evalúan diferentes niveles de temperatura y humedad relativa.

Este diseño permite aplicar análisis de varianza (ANOVA) para determinar si existen diferencias significativas entre los tratamientos evaluados.

\subsection{Análisis estadístico}

El análisis de varianza (ANOVA) permite determinar si las variables de proceso generan diferencias significativas en la calidad del café.

En caso de encontrar diferencias significativas, se aplica la prueba de Tukey para identificar específicamente entre qué tratamientos existen dichas diferencias.

Este enfoque permite validar experimentalmente los resultados obtenidos.

\subsection{Modelo de regresión}

Se propone un modelo de regresión lineal múltiple para cuantificar la relación entre las variables de proceso y las variables de respuesta:

\[
y = \beta_0 + \beta_1 T + \beta_2 HR + \epsilon
\]

donde:

\begin{itemize}
    \item $T$: temperatura
    \item $HR$: humedad relativa
    \item $y$: variable de respuesta
\end{itemize}

\subsection{Pipeline de procesamiento de datos}

En escenarios reales, los datos pueden presentar valores faltantes o inconsistencias debido a fallas en sensores o condiciones experimentales.

Por esta razón, se propone el uso de un pipeline que integre:

\begin{itemize}
    \item Imputación de datos faltantes
    \item Escalado de variables
    \item Modelo de regresión
\end{itemize}

Esto permite automatizar el análisis y garantizar la robustez del modelo.

\subsection{Ejemplo de implementación}

\begin{lstlisting}
import pandas as pd
from sklearn.pipeline import Pipeline
from sklearn.impute import SimpleImputer
from sklearn.preprocessing import StandardScaler
from sklearn.linear_model import LinearRegression

# Datos simulados
df = pd.DataFrame({
    "Temperatura": [40, 40, 50, 50],
    "Humedad": [50, 60, 50, 60],
    "Calidad": [80, 75, 85, 78]
})

X = df[["Temperatura", "Humedad"]]
y = df["Calidad"]

pipeline = Pipeline([
    ("imputador", SimpleImputer(strategy="median")),
    ("escalador", StandardScaler()),
    ("modelo", LinearRegression())
])

pipeline.fit(X, y)

pred = pipeline.predict([[45, 55]])
print(pred)
\end{lstlisting}

\subsection{Integración metodológica}

El enfoque propuesto combina herramientas estadísticas y de aprendizaje automático:

\begin{itemize}
    \item ANOVA: Validación de diferencias entre tratamientos
    \item Tukey: Comparación entre grupos específicos
    \item Regresión: Modelado de relaciones entre variables
    \item Pipeline: Automatización del procesamiento de datos
\end{itemize}

Esta integración permite desarrollar modelos robustos que no solo explican el fenómeno, sino que también permiten realizar predicciones en condiciones reales.